\documentclass[12pt,a4paper]{article}

\usepackage[T1]{fontenc}
\usepackage[utf8]{inputenc}
\usepackage[brazilian]{babel}
\usepackage{lmodern}
\usepackage[a4paper,left=3cm,right=2cm,top=3cm,bottom=2cm]{geometry}
\usepackage{setspace}
\usepackage{booktabs}
\usepackage{array}
\usepackage{enumitem}
\usepackage{xcolor}
\usepackage{tikz}
\usepackage[most]{tcolorbox}
\usepackage[hidelinks]{hyperref}

\setstretch{1.25}
\setlength{\parindent}{1.25cm}
\setlist{itemsep=2pt,topsep=4pt}
\usepackage{pdflscape}
% ===============================================================
% Material desplugado: ícones, peças, barco, tabuleiro (TikZ)
% ===============================================================
\usepackage{graphicx}
\usetikzlibrary{decorations.pathmorphing,shapes.geometric,arrows.meta}

% ---------- Ícones (desenhados em torno da origem, ~2 x 2 unidades) ----------
\newcommand{\iconeFazendeiro}{%
  \fill[red!65!black!70] (-0.72,-0.95) -- (-0.62,-0.42) -- (0.62,-0.42) -- (0.72,-0.95) -- cycle;
  \fill[blue!65!black] (-0.38,-0.95) -- (-0.34,-0.42) -- (0.34,-0.42) -- (0.38,-0.95) -- cycle;
  \fill[orange!35!pink] (0,0) circle (0.5);
  \fill[black] (-0.18,0.0) circle (0.05);
  \fill[black] (0.18,0.0) circle (0.05);
  \draw[thick,line cap=round] (-0.17,-0.2) .. controls (-0.07,-0.3) and (0.07,-0.3) .. (0.17,-0.2);
  \fill[yellow!75!orange] (-0.42,0.36) .. controls (-0.42,1.0) and (0.42,1.0) .. (0.42,0.36) -- cycle;
  \fill[yellow!75!orange] (0,0.36) ellipse (0.85 and 0.16);
  \draw[orange!50!black,thin] (0,0.36) ellipse (0.85 and 0.16);
  \fill[red!70!black] (-0.42,0.42) rectangle (0.42,0.56);
}

\newcommand{\iconeLobo}{%
  \draw[fill=gray!55,draw=black!70,thick,line join=round]
    (-0.78,-0.05) -- (-0.88,0.98) -- (-0.28,0.58) -- (0.28,0.58) -- (0.88,0.98)
    -- (0.78,-0.05) -- (0.4,-0.78) -- (-0.4,-0.78) -- cycle;
  \fill[pink!70] (-0.74,0.72) -- (-0.5,0.55) -- (-0.72,0.32) -- cycle;
  \fill[pink!70] (0.74,0.72) -- (0.5,0.55) -- (0.72,0.32) -- cycle;
  \fill[gray!15] (0,-0.42) ellipse (0.42 and 0.33);
  \fill[yellow!85!orange] (-0.34,0.2) ellipse (0.14 and 0.08);
  \fill[yellow!85!orange] (0.34,0.2) ellipse (0.14 and 0.08);
  \fill[black] (-0.34,0.2) circle (0.04);
  \fill[black] (0.34,0.2) circle (0.04);
  \fill[black] (0,-0.2) ellipse (0.11 and 0.08);
  \draw[black,thick] (0,-0.27) -- (0,-0.42);
  \draw[black,thick] (-0.22,-0.5) .. controls (-0.1,-0.42) .. (0,-0.42)
    .. controls (0.1,-0.42) .. (0.22,-0.5);
  \fill[white] (-0.17,-0.46) -- (-0.1,-0.63) -- (-0.03,-0.44) -- cycle;
  \fill[white] (0.17,-0.46) -- (0.1,-0.63) -- (0.03,-0.44) -- cycle;
}

\newcommand{\iconeOvelha}{%
  \foreach \x/\y in {-0.55/0.25, 0/0.5, 0.55/0.25, -0.65/-0.2, 0.65/-0.2, -0.35/-0.55, 0.35/-0.55}
    \draw[fill=white,draw=gray!70,thick] (\x,\y) circle (0.4);
  \fill[white] (0,0) circle (0.62);
  \fill[gray!45!black] (-0.45,-0.02) ellipse (0.19 and 0.09);
  \fill[gray!45!black] (0.45,-0.02) ellipse (0.19 and 0.09);
  \fill[gray!45!black] (0,-0.22) ellipse (0.3 and 0.4);
  \fill[white] (-0.12,-0.12) circle (0.05);
  \fill[white] (0.12,-0.12) circle (0.05);
  \fill[pink!60] (0,-0.46) ellipse (0.1 and 0.06);
}

\newcommand{\iconeAlface}{%
  \foreach \a in {0,72,144,216,288}{%
    \begin{scope}[rotate=\a]
      \draw[fill=green!55!black!85,draw=green!25!black,thick] (0,0.12) ellipse (0.34 and 0.8);
    \end{scope}}
  \foreach \a in {36,108,180,252,324}{%
    \begin{scope}[rotate=\a]
      \draw[fill=lime!70!green,draw=green!30!black] (0,0.06) ellipse (0.27 and 0.56);
    \end{scope}}
  \fill[lime!45!yellow] (0,0) circle (0.22);
}

\newcommand{\iconeCachorro}{%
  \fill[brown!55!black,rotate around={18:(-0.7,0.0)}] (-0.7,0.0) ellipse (0.22 and 0.55);
  \fill[brown!55!black,rotate around={-18:(0.7,0.0)}] (0.7,0.0) ellipse (0.22 and 0.55);
  \draw[fill=brown!55!orange,draw=brown!40!black,thick] (0,0.0) ellipse (0.62 and 0.66);
  \fill[brown!35!black] (0.28,0.32) ellipse (0.18 and 0.2);
  \fill[brown!15!white] (0,-0.28) ellipse (0.36 and 0.28);
  \fill[red!70] (0,-0.48) ellipse (0.08 and 0.1);
  \fill[black] (0,-0.15) ellipse (0.13 and 0.09);
  \draw[black,thick] (0,-0.24) -- (0,-0.34);
  \draw[black,thick] (-0.14,-0.4) .. controls (-0.06,-0.34) .. (0,-0.34)
    .. controls (0.06,-0.34) .. (0.14,-0.4);
  \fill[black] (-0.25,0.15) circle (0.07);
  \fill[black] (0.25,0.15) circle (0.07);
  \fill[white] (-0.27,0.18) circle (0.02);
  \fill[white] (0.23,0.18) circle (0.02);
}

\newcommand{\iconeGalinha}{%
  \draw[orange!80!black,line width=1.5pt] (-0.22,-0.72) -- (-0.22,-0.95) -- (-0.38,-0.97);
  \draw[orange!80!black,line width=1.5pt] (0.22,-0.72) -- (0.22,-0.95) -- (0.38,-0.97);
  \draw[fill=white,draw=gray!60,thick] (0,-0.1) circle (0.72);
  \draw[fill=gray!15,draw=gray!50] (-0.6,-0.2) ellipse (0.2 and 0.38);
  \draw[fill=gray!15,draw=gray!50] (0.6,-0.2) ellipse (0.2 and 0.38);
  \fill[red!75] (-0.18,0.6) circle (0.13);
  \fill[red!75] (0,0.7) circle (0.15);
  \fill[red!75] (0.18,0.6) circle (0.13);
  \fill[black] (-0.28,0.3) circle (0.06);
  \fill[black] (0.28,0.3) circle (0.06);
  \fill[red!75] (0,-0.2) ellipse (0.07 and 0.1);
  \fill[orange] (-0.15,0.15) -- (0.15,0.15) -- (0,-0.1) -- cycle;
}

% ---------- Peças (comandos de desenho; origem no canto inferior esquerdo) ----------
% Peça quadrada (peso 1): cor, ícone, nome, peso
\newcommand{\pecaS}[4]{%
  \draw[rounded corners=4pt,line width=1pt,draw=#1!70!black,fill=#1!12] (0,0) rectangle (2.6,2.6);
  \begin{scope}[shift={(1.3,1.38)},scale=0.82] #2 \end{scope}
  \node[font=\bfseries\scriptsize,text=#1!50!black] at (1.3,0.24) {#3};
  \node[draw=#1!70!black,fill=white,rounded corners=2pt,inner sep=2pt,font=\bfseries\scriptsize,text=#1!50!black,anchor=north west] at (0.1,2.5) {peso #4};
}

% Peça larga (peso 2): cor, ícone, nome, descrição
\newcommand{\pecaL}[4]{%
  \draw[rounded corners=4pt,line width=1pt,draw=#1!70!black,fill=#1!12] (0,0) rectangle (5.4,2.6);
  \begin{scope}[shift={(1.3,1.38)},scale=0.82] #2 \end{scope}
  \node[font=\bfseries\large,text=#1!50!black,anchor=west] at (2.75,1.65) {#3};
  \node[font=\small,text=#1!50!black,anchor=west] at (2.75,1.0) {#4};
  \node[draw=#1!70!black,fill=white,rounded corners=2pt,inner sep=2pt,font=\bfseries\scriptsize,text=#1!50!black,anchor=north west] at (0.1,2.5) {peso 2};
  \node[font=\scriptsize,text=#1!50!black,anchor=west] at (2.75,0.4) {ocupa 2 espaços};
}

% Barco (8,6 x 3,4 cm): três espaços de 2,8 cm; o primeiro traz o Fazendeiro
\newcommand{\barcoCmd}{%
  \draw[rounded corners=8pt,line width=1.2pt,draw=brown!60!black,fill=brown!30] (0,0) rectangle (8.6,3.4);
  \draw[fill=white!85!brown,draw=brown!60!black,thick] (0.1,0.62) rectangle (2.9,3.3);
  \draw[dashed,draw=brown!60!black,thick,fill=white] (2.9,0.62) rectangle (8.5,3.3);
  \draw[dotted,brown!60!black,thick] (5.7,0.62) -- (5.7,3.3);
  \begin{scope}[shift={(1.5,2.1)},scale=0.85] \iconeFazendeiro \end{scope}
  \node[font=\bfseries\scriptsize,text=brown!50!black] at (1.5,1.08) {Fazendeiro};
  \node[font=\bfseries\scriptsize,text=brown!50!black] at (1.5,0.80) {peso 1};
  \node[font=\scriptsize,text=brown!50!black] at (4.3,2.0) {espaço de carga};
  \node[font=\scriptsize,text=brown!50!black] at (4.3,1.6) {(cabe peso 1)};
  \node[font=\scriptsize,text=brown!50!black] at (7.1,2.0) {espaço de carga};
  \node[font=\scriptsize,text=brown!50!black] at (7.1,1.6) {(cabe peso 1)};
  \node[font=\bfseries\small,text=brown!40!black] at (4.3,0.31) {BARCO: peso máximo 3};
}

% Kit completo de um grupo (origem no canto superior esquerdo, 8,6 x 12,1 cm)
\newcommand{\kitCmd}{%
  \begin{scope}[shift={(0,-2.6)}]\pecaL{gray}{\iconeLobo}{Lobo}{Predador}\end{scope}
  \begin{scope}[shift={(5.7,-2.6)}]\pecaS{green}{\iconeAlface}{Alface}{1}\end{scope}
  \begin{scope}[shift={(0,-5.4)}]\pecaL{blue}{\iconeOvelha}{Ovelha}{Herbívora}\end{scope}
  \begin{scope}[shift={(5.7,-5.4)}]\pecaS{brown}{\iconeCachorro}{Cachorro}{1}\end{scope}
  \begin{scope}[shift={(0,-8.2)}]\pecaS{orange}{\iconeGalinha}{Galinha}{1}\end{scope}
  \node[font=\scriptsize,text width=5.3cm,align=left,anchor=west] at (3.0,-6.9)
    {Recortar pelo contorno. Peça larga = peso 2; peça quadrada = peso 1. O Fazendeiro já está impresso no barco.};
  \begin{scope}[shift={(0,-12.1)}]\barcoCmd\end{scope}
}

% Folha de peças: quatro kits por página (A4 retrato, margens de 1,5 cm)
\newcommand{\paginapecas}{%
  \begin{tikzpicture}
    \begin{scope}[shift={(0,0)}]\kitCmd\end{scope}
    \begin{scope}[shift={(9.2,0)}]\kitCmd\end{scope}
    \begin{scope}[shift={(0,-13.4)}]\kitCmd\end{scope}
    \begin{scope}[shift={(9.2,-13.4)}]\kitCmd\end{scope}
  \end{tikzpicture}}

% ---------- Gota de combustível ----------
\newcommand{\gota}[4]{% x, y, cor, rótulo
  \begin{scope}[shift={(#1,#2)}]
    \draw[thick,fill=#3] (0,0.62) .. controls (0.12,0.35) and (0.45,0.12) .. (0.45,-0.22)
      arc (0:-180:0.45) .. controls (-0.45,0.12) and (-0.12,0.35) .. (0,0.62) -- cycle;
    \node[font=\large\bfseries] at (0,-0.2) {#4};
  \end{scope}}

% ---------- Tabuleiro (27,7 x 19 cm, para A4 paisagem com margens de 0,9 cm) ----------
\newcommand{\linharegra}[6]{% y, id, ícone A, ícone B, texto, largura do texto
  \node[font=\large\bfseries,text=violet!70!black] at (10.05,#1) {#2};
  \begin{scope}[shift={(11.2,#1)},scale=0.5]#3\end{scope}
  \node[font=\Large\bfseries] at (12.1,#1) {$\times$};
  \begin{scope}[shift={(13.0,#1)},scale=0.5]#4\end{scope}
  \node[anchor=west,text width=#6,font=\footnotesize] at (13.75,#1) {#5};
}

\newcommand{\tabuleiro}{%
\begin{tikzpicture}[x=1cm,y=1cm]
  % fundos
  \fill[green!12] (0,2.8) rectangle (9.2,17.2);
  \fill[green!12] (18.5,2.8) rectangle (27.7,17.2);
  \fill[cyan!18]  (9.2,2.8) rectangle (18.5,17.2);
  \draw[cyan!50!blue!60,decorate,decoration={snake,amplitude=0.6mm,segment length=6mm},thick] (9.4,3.0) -- (9.4,17.0);
  \draw[cyan!50!blue!60,decorate,decoration={snake,amplitude=0.6mm,segment length=6mm},thick] (18.3,3.0) -- (18.3,17.0);
  \draw[thick,black!60] (0,2.8) rectangle (27.7,17.2);
  \draw[thick,black!60] (9.2,2.8) -- (9.2,17.2);
  \draw[thick,black!60] (18.5,2.8) -- (18.5,17.2);
  % título
  \fill[blue!60!black] (0,17.4) rectangle (27.7,19);
  \node[white,font=\Huge\bfseries] at (13.85,18.2) {TRAVESSIA DO RIO};
  % margens
  \node[font=\Large\bfseries,text=green!40!black] at (4.6,16.65) {MARGEM ESQUERDA (E)};
  \node[font=\Large\bfseries,text=green!40!black] at (23.1,16.65) {MARGEM DIREITA (D)};
  \foreach \c in {0,1,2}{\foreach \r in {0,1,2}{
    \draw[dashed,green!40!black,thick] ({0.4+2.8*\c},{16.1-2.8*(\r+1)}) rectangle ++(2.8,2.8);
    \draw[dashed,green!40!black,thick] ({18.9+2.8*\c},{16.1-2.8*(\r+1)}) rectangle ++(2.8,2.8);
  }}
  % cais do barco
  \draw[dashed,brown!60!black,very thick] (0.3,3.6) rectangle (8.9,7.0);
  \draw[dashed,brown!60!black,very thick] (18.8,3.6) rectangle (27.4,7.0);
  \node[font=\small\bfseries,text=brown!50!black] at (4.6,3.2) {CAIS E: posição do barco};
  \node[font=\small\bfseries,text=brown!50!black] at (23.1,3.2) {CAIS D: posição do barco};
  % painel de regras
  \draw[rounded corners=8pt,fill=white,draw=black!50,thick] (9.6,3.4) rectangle (18.1,16.8);
  \node[font=\large\bfseries] at (13.85,16.2) {REGRAS DE CONFLITO};
  \node[font=\scriptsize,text width=8.0cm,align=center,anchor=north] at (13.85,15.8)
    {Valem para qualquer par de pe\c{c}as na margem em que o Fazendeiro (isto \'e, o barco) \emph{n\~ao} est\'a, mesmo com outras pe\c{c}as junto.};
  \linharegra{13.3}{R1}{\iconeLobo}{\iconeOvelha}{Lobo devora Ovelha, \textbf{exceto} se o Cachorro estiver na margem.}{3.15cm}
  \begin{scope}[shift={(17.65,13.3)},scale=0.4]\iconeCachorro\end{scope}
  \draw[gray!40] (9.9,12.1) -- (17.8,12.1);
  \linharegra{10.9}{R2}{\iconeOvelha}{\iconeAlface}{Ovelha devora Alface.}{3.9cm}
  \draw[gray!40] (9.9,9.7) -- (17.8,9.7);
  \linharegra{8.5}{R3}{\iconeGalinha}{\iconeAlface}{Galinha bica Alface.}{3.9cm}
  \draw[gray!40] (9.9,7.3) -- (17.8,7.3);
  \linharegra{6.1}{R4}{\iconeCachorro}{\iconeGalinha}{Cachorro ataca Galinha.}{3.9cm}
  \node[font=\scriptsize,text width=7.8cm,align=center] at (13.85,4.3)
    {Pe\c{c}a larga = peso 2 $\cdot$ pe\c{c}a quadrada = peso 1\\
     Barco: peso m\'aximo 3 (Fazendeiro, peso 1, mais carga de at\'e peso 2)};
  % combustível
  \fill[yellow!15] (0,0) rectangle (27.7,2.6);
  \draw[thick,black!60] (0,0) rectangle (27.7,2.6);
  \node[anchor=west,font=\small\bfseries] at (0.3,2.2)
    {TANQUE DE COMBUSTÍVEL: risque uma gota a cada travessia, mesmo as que forem desfeitas};
  \node[anchor=east,font=\scriptsize] at (27.4,2.2)
    {Rodada 1: use as gotas 1 a 7. Rodada 2: use as 13 gotas.};
  \foreach \i in {1,...,13}{
    \gota{{1.4+2.0*(\i-1)}}{0.85}{orange!20}{\i}
  }
  \gota{21.4}{0.85}{yellow!60}{11}
  \gota{25.4}{0.85}{red!20}{13}
  \node[font=\scriptsize\bfseries,text=black!70] at (1.4,1.78) {tanque cheio};
  \node[font=\scriptsize\bfseries,text=orange!60!black] at (21.4,1.78) {consumo mínimo};
  \node[font=\scriptsize\bfseries,text=red!70!black] at (25.4,1.78) {tanque vazio};
\end{tikzpicture}}

% ---------- Carta de dica (5,0 x 3,9 cm) ----------
\newcommand{\dica}[4][2]{% [rodada] número, desconto na nota, texto
  \begin{tikzpicture}[baseline=0]
    \draw[rounded corners=6pt,thick,fill=yellow!10,draw=black!70] (0,0) rectangle (5.0,5.2);
    \node[anchor=north west,font=\bfseries\large] at (0.2,5.05) {Dica #2};
    \node[anchor=north east,font=\scriptsize] at (4.8,5.0) {rodada #1};
    \node[anchor=north west,font=\scriptsize\itshape] at (0.2,4.52) {desconta #3 da nota};
    \node[anchor=north west,text width=4.5cm,font=\small,align=left] at (0.25,4.05) {#4};
  \end{tikzpicture}}


% ---------------------------------------------------------------
% Macro para cartas impressas (5,0 cm x 6,4 cm)
% #1 cor | #2 título | #3 subtítulo | #4 corpo
% ---------------------------------------------------------------
\newcommand{\carta}[4]{%
\begin{minipage}[t]{5.0cm}
\begin{tcolorbox}[enhanced, colback=white, colframe=#1, colbacktitle=#1,
  coltitle=white, title={\bfseries #2}, halign title=center, boxrule=1pt,
  arc=2.5mm, height=6.4cm, valign=top, left=2mm, right=2mm, top=1.5mm,
  bottom=1mm, fontupper=\small, halign=flush left]
\textbf{#3}\par\smallskip
#4
\end{tcolorbox}
\end{minipage}}

% ---------- Fluxograma da rotina de verificação ----------
\newcommand{\lacunaS}{\rule{2.4cm}{0.4pt}\\[0.45cm]\rule{2.4cm}{0.4pt}\\[0.45cm]\rule{2.4cm}{0.4pt}}
\newcommand{\lacunaL}{\rule{3.8cm}{0.4pt}\\[0.45cm]\rule{3.8cm}{0.4pt}}

\newcommand{\fluxograma}[4]{% decisão da carga, decisão da regra, ação de desfazer, decisão final
\begin{tikzpicture}[
  caixa/.style={draw,rounded corners=4pt,thick,fill=blue!6,text width=4.4cm,align=center,minimum height=1.0cm,font=\small},
  decisao/.style={draw,diamond,aspect=2.0,thick,fill=yellow!14,text width=3.0cm,align=center,inner sep=1pt,font=\small},
  seta/.style={-{Stealth[length=2.2mm]},thick},
  rot/.style={font=\footnotesize\bfseries,fill=white,inner sep=1pt}]
  \node[caixa,fill=green!14] (ini) at (3,0) {Início: peças na margem E, tanque cheio};
  \node[caixa] (esc) at (3,-1.7) {Escolher a carga};
  \node[decisao] (d1) at (3,-4.3) {#1};
  \node[caixa] (atr) at (3,-7.5) {Atravessar e riscar 1 gota de combustível (tanque vazio: recomeçar)};
  \node[decisao] (d2) at (3,-10.7) {#2};
  \node[caixa] (des) at (10.4,-10.7) {#3};
  \node[caixa] (reg) at (3,-14.0) {Registrar a travessia na Ficha};
  \node[decisao] (d3) at (3,-17.0) {#4};
  \node[caixa,fill=green!14] (fim) at (3,-20.0) {Fim};
  \draw[seta] (ini) -- (esc);
  \draw[seta] (esc) -- (d1);
  \draw[seta] (d1.south) -- node[rot,right=2pt]{Sim} (atr.north);
  \draw[seta] (atr) -- (d2);
  \draw[seta] (d2.east) -- node[rot,above]{Sim} (des.west);
  \draw[seta] (d2.south) -- node[rot,right=2pt]{Não} (reg.north);
  \draw[seta] (reg) -- (d3);
  \draw[seta] (d3.south) -- node[rot,right=2pt]{Sim} (fim.north);
  \draw[thick] (d1.east) -- node[rot,above]{Não} (13.6,-4.3);
  \draw[thick] (d3.east) -- node[rot,above]{Não} (13.6,-17.0);
  \draw[thick] (des.east) -- (13.6,-10.7);
  \draw[thick] (13.6,-17.0) -- (13.6,-4.3);
  \draw[seta] (13.6,-4.3) |- (esc.east);
\end{tikzpicture}}

\title{\textbf{Plano de Aula}\\[4pt]
\large Travessia do Rio: um jogo de cartas desplugado\\
para trabalhar a ideia de algoritmo com Aprendizagem Baseada em Problemas}
\author{}
\date{}

\begin{document}
\maketitle
\vspace{-1.5cm}

% ===============================================================
\section{Identificação}

\begin{tabular}{@{}>{\bfseries}p{4.9cm}p{9.8cm}@{}}
Componente curricular & Educação em Computação \\
Instituição/Programa  & \rule[-2pt]{7cm}{0.4pt} \\
Docente(s)            & \rule[-2pt]{7cm}{0.4pt} \\
Duração               & 45 minutos \\
Público               & Estudantes de mestrado de áreas diversas; não se exige conhecimento prévio de Computação \\
Organização           & Grupos de três a quatro estudantes, de preferência com formações diferentes \\
Metodologia           & Aprendizagem Baseada em Problemas, com um único problema de curta duração \\
Recurso               & Atividade desplugada: tabuleiro, peças e fluxograma em papel \\
Tema                  & Algoritmo como passo a passo: repetição, decisão e verificação \\
\end{tabular}

% ===============================================================
\section{Justificativa}

O problema da travessia do rio com o lobo, a ovelha e a alface é um enigma
antigo, tradicionalmente atribuído a Alcuíno de York. Na forma
original, a solução aparece rápido por tentativa e erro e tem 7 travessias.
Esta proposta usa o enigma como pretexto para discutir o que é um algoritmo, ou
seja, um passo a passo que qualquer pessoa consegue executar sem precisar
interpretá-lo. Para isso o jogo foi ampliado: são cinco personagens com pesos
diferentes, um barco com apenas dois espaços de carga, uma regra com exceção (o
Cachorro protege a Ovelha do Lobo, mas ameaça a Galinha) e um tanque de
combustível limitado. A melhor solução tem 11 travessias, valor confirmado por
programa de computador antes da elaboração deste plano.

A turma reúne estudantes de várias áreas, e por isso a atividade não depende de
conteúdos de Computação. Pede-se aos grupos que descrevam com clareza o que
fazem a cada travessia, percebam o que se repete e onde há decisões do tipo
``se isto, então aquilo'', e representem essa rotina em um fluxograma simples.
A atividade é desplugada, no sentido de \textit{Computer Science Unplugged}
\cite{bell2009}. A aula segue a Aprendizagem Baseada em Problemas: o
problema é apresentado antes de qualquer explicação, os estudantes trabalham em
pequenos grupos, definem por conta própria como resolvê-lo e o docente atua
como facilitador, e não como transmissor de conteúdo \cite{barrows1986,hmelo2004},
no quadro das metodologias ativas \cite{bacich2018}. Dentro dos 40 minutos, o
problema é deliberadamente fechado e com regras explícitas; o caráter aberto
de um problema real é atendido apenas em parte.

% ===============================================================
\section{Objetivos de aprendizagem}

Os objetivos usam verbos da taxonomia de Bloom revisada \cite{anderson2001}.
Ao final da aula, espera-se que o estudante seja capaz de:

\begin{enumerate}[label=\arabic*.]
  \item descrever, em linguagem simples, um passo a passo para resolver um
        problema com regras;
  \item identificar, nesse passo a passo, o que se repete e onde há decisões
        do tipo ``se\ldots então'';
  \item verificar cada ação contra as regras e corrigir o rumo quando houver erro;
  \item representar o passo a passo em um fluxograma simples e a situação do
        jogo em um diagrama de duas margens;
  \item discutir o uso do jogo como recurso didático em turmas de formações
        diferentes.
\end{enumerate}

% ===============================================================
\section{Problema norteador e produto}

\textbf{Problema:} como escrever um passo a passo que qualquer pessoa
consiga seguir para levar todos à outra margem sem quebrar nenhuma regra e
gastando pouco combustível?

\textbf{Produto:} a Ficha de Registro do grupo (Apêndice B), com as travessias
anotadas, o fluxograma da rotina de verificação preenchido e as respostas às
questões.

% ===============================================================
\section{O jogo}

\subsection{Componentes}

O material completo, com ilustrações originais, está no Apêndice~A, pronto para
imprimir e recortar. Cada grupo recebe:

\begin{itemize}
  \item 1 tabuleiro em A4 paisagem (Figura~\ref{fig:tabuleiro}), com as duas
        margens (nove espaços cada), os cais do barco, as regras de conflito
        R1 a R4 e o tanque de combustível de 13 gotas;
  \item 5 peças de personagem (Figura~\ref{fig:pecas}): Lobo e Ovelha são peças
        largas, de peso 2; Alface, Cachorro e Galinha são peças quadradas, de
        peso 1;
  \item 1 barco que aguenta peso máximo 3, com três espaços de peso 1. O
        primeiro é ocupado pelo Fazendeiro, que nunca sai do barco;
  \item 1 regulamento de uma página e a Ficha de Registro;
  \item 8 cartas de dica (3 para a rodada 1 e 5 para a rodada 2), guardadas pelo docente e entregues quando pedidas.
\end{itemize}

O peso de cada peça aparece no tamanho dela e em uma etiqueta numerada. Como o
barco tem só dois espaços de carga, o peso máximo se confere olhando, sem fazer
conta.

O material usa três palavras distintas, e vale explicitá-las aos grupos. O
\textit{peso} diz quanto a peça ocupa no barco. A \textit{gota} mede o
combustível gasto, uma por travessia. Os \textit{pontos} só existem na
\textit{nota} da avaliação, atribuída pelo docente. O jogo em si não tem
pontuação.

\begin{figure}[ht]
\centering
\resizebox{\textwidth}{!}{\tabuleiro}
\caption{Tabuleiro do jogo, reduzido. A versão em tamanho real está no
Apêndice~A.}
\label{fig:tabuleiro}
\end{figure}

\begin{figure}[ht]
\centering
\resizebox{0.5\textwidth}{!}{\begin{tikzpicture}\kitCmd\end{tikzpicture}}
\caption{Peças e barco de um grupo, reduzidos.}
\label{fig:pecas}
\end{figure}

\subsection{Papéis no grupo}

Em cada rodada, os integrantes se dividem em papéis, que trocam de mãos na
rodada seguinte:

\begin{itemize}
  \item o \textit{operador} move as peças e o barco;
  \item o \textit{fiscal} lê as regras R1 a R4 e confere cada travessia antes de
        ela ser aceita;
  \item o \textit{anotador} preenche a Ficha;
  \item o \textit{guardião do combustível} risca uma gota a cada travessia; em
        grupos de três, o anotador também faz isso.
\end{itemize}

O fiscal faz, na prática, o trabalho de validação que o fluxograma vai
descrever depois. Dividir assim o trabalho dá a todos algo a fazer, qualquer que
seja a formação.

\subsection{Regras}

\begin{enumerate}[label=\arabic*.]
  \item \textbf{Objetivo.} Levar o Lobo, a Ovelha, a Alface, o Cachorro e a
        Galinha da margem E para a margem D sem quebrar nenhuma regra de
        conflito e sem esvaziar o tanque.
  \item \textbf{Rodadas.} A rodada 1 é um aquecimento: o Lobo fica fora do jogo e
        o tanque tem 7 gotas. Na rodada 2 entram todos os personagens e o
        tanque tem 13 gotas.
  \item \textbf{Preparação.} As peças começam nos espaços da margem E e o barco,
        no cais E. As regras R1 a R4 estão impressas no centro do tabuleiro.
  \item \textbf{Travessia.} Só o Fazendeiro conduz o barco, e ele nunca o deixa.
        O barco aguenta peso total 3, e o Fazendeiro pesa 1. Sobra peso 2 para
        a carga: uma peça larga (peso 2: Lobo ou Ovelha), até duas peças
        quadradas (peso 1 cada: Alface, Cachorro, Galinha) ou nenhuma carga. Ao
        atravessar, o barco passa de um cais ao outro e a carga desembarca na
        margem de chegada.
  \item \textbf{Verificação.} Logo após cada travessia, o fiscal examina a
        margem em que o barco \emph{não} está. Se alguma regra se aplicar a
        ela, a travessia é inválida e deve ser desfeita. Na margem em que o
        barco está, e dentro dele, não há conflito. As regras valem para
        qualquer par de peças dessa margem, mesmo que haja outras peças junto.
        Só o Fazendeiro evita o conflito, e a única exceção é o Cachorro, que
        protege a Ovelha do Lobo (R1).
  \item \textbf{Combustível.} Cada travessia gasta uma gota, inclusive as que
        forem desfeitas. Se o tanque esvaziar antes de todos chegarem à margem
        D, o grupo recomeça a rodada e anota a tentativa anterior no verso da
        Ficha.
  \item \textbf{Planejamento.} Simular as travessias no papel é permitido e não
        gasta combustível.
  \item \textbf{Dicas.} Há dicas para cada rodada (três na rodada 1 e cinco na
        rodada 2), que podem ser pedidas ao docente. Cada uma reduz a nota do
        grupo: 0,25, 0,50 e 0,75 ponto na rodada 1; 0,25, 0,25, 0,50, 0,50 e
        0,75 ponto na rodada 2, na ordem das cartas.
\end{enumerate}

% ===============================================================
\section{Cronograma da aula (40 minutos)}

\begin{center}
\renewcommand{\arraystretch}{1.1}
\small
\begin{tabular}{@{}>{\raggedright\arraybackslash}p{2.2cm}>{\raggedright\arraybackslash}p{3.6cm}>{\raggedright\arraybackslash}p{8.7cm}@{}}
\toprule
\textbf{Tempo} & \textbf{Fase} & \textbf{Atividades} \\
\midrule
0 a 5 min   & Apresentação      & Apresentação do problema, formação dos grupos e dos papéis, leitura do regulamento e demonstração de uma travessia no quadro, com a verificação feita em voz alta. \\
5 a 11 min  & Rodada 1          & Aquecimento sem o Lobo, com tanque de 7 gotas. Os grupos conhecem as peças e a rotina de verificação. \\
11 a 26 min & Rodada 2          & Jogo completo, com tanque de 13 gotas. Registro das travessias na Ficha. \\
26 a 34 min & Fluxograma        & Parte 2 da Ficha: preencher o fluxograma da rotina de verificação. Parte 3: responder à Q1. \\
34 a 40 min & Socialização      & Dois grupos apresentam seus fluxogramas; comparação (Q2) e fechamento (Q3, que pode ser concluída em casa). \\
\bottomrule
\end{tabular}
\end{center}

A divisão de tempo é uma estimativa e não foi testada com estudantes. Convém
fazer um pré-teste com um grupo pequeno antes da aplicação, principalmente da
rodada 2.

% ===============================================================
\section{Condução da aula}

\begin{itemize}
  \item Na demonstração inicial, execute uma travessia no quadro dizendo em voz
        alta cada verificação. É essa fala que os grupos vão depois transformar
        em fluxograma.
  \item Forme grupos mistos, com pelo menos uma pessoa menos familiarizada com
        Computação em cada um, e peça que o papel de fiscal passe por todos.
  \item A estratégia da versão clássica (levar primeiro a Ovelha) quebra uma
        regra. Na rodada 2 existe um único primeiro movimento válido, e descobri-lo
        é uma boa discussão sobre verificar antes de agir.
  \item A solução exige uma viagem sem carga e voltas com carga. A Alface, por
        exemplo, cruza o rio cinco vezes. Vale comentar que andar para trás
        também faz parte do plano.
  \item Quando um grupo terminar cedo, ofereça as cartas de perturbação da
        seção de extensões.
  \item Antes da Parte 2, mostre em um minuto os dois símbolos do fluxograma: o
        retângulo é uma ação e o losango é uma pergunta de sim ou não.
  \item Na socialização, projete ou distribua o fluxograma de referência do
        Apêndice~C e a ficha preenchida do Apêndice~D só depois de os grupos
        apresentarem os seus, para que a comparação parta do que eles mesmos
        construíram.
\end{itemize}

% ===============================================================
\section{Ideias de computação trabalhadas}

\begin{itemize}
  \item \textit{Algoritmo:} um passo a passo claro o bastante para que outra
        pessoa o execute sem adivinhar nada \cite{wing2006}.
  \item \textit{Repetição:} escolher a carga, atravessar e conferir se repetem
        até que todos estejam na margem D.
  \item \textit{Condição:} cada verificação é uma pergunta de sim ou não que
        leva a caminhos diferentes.
  \item \textit{Validação:} conferir cada ação contra as regras antes de
        aceitá-la.
  \item \textit{Depuração:} desfazer uma travessia inválida, achar o motivo e
        tentar de outro jeito.
  \item \textit{Diagrama e fluxograma:} o quadro das duas margens mostra a
        situação do jogo, e o fluxograma mostra a rotina.
  \item \textit{Otimização simples:} procurar o menor consumo de combustível.
\end{itemize}

% ===============================================================
\section{Avaliação}

\begin{center}
\renewcommand{\arraystretch}{1.3}
\begin{tabular}{@{}>{\raggedright\arraybackslash}p{8.2cm}c>{\raggedright\arraybackslash}p{4.6cm}@{}}
\toprule
\textbf{Critério} & \textbf{Pontos da nota} & \textbf{Evidência na Ficha} \\
\midrule
Solução válida da rodada 2 (3,0 pontos com 11 travessias válidas; 2,0 pontos com 13) & 3,0 & Parte 1 \\
Fluxograma completo e coerente com o que o grupo fez & 3,0 & Parte 2 \\
Identificação da repetição e das verificações no fluxograma & 2,0 & Parte 2 \\
Análise de erros e reflexão didática & 2,0 & Q1 a Q3 \\
\midrule
\textbf{Total} & \textbf{10,0} & \\
\bottomrule
\end{tabular}
\end{center}

Descontos por dica: 0,25, 0,50 e 0,75 ponto na rodada 1; 0,25, 0,25, 0,50, 0,50 e 0,75 ponto na rodada 2, conforme a regra 8.

% ===============================================================
\section{Extensões}

\subsection{Cartas de perturbação}

Para completar o ciclo de resolução do problema fora da sala, cada grupo escolhe uma
alteração das regras, prevê o efeito e depois confere. Os resultados abaixo
foram obtidos com o programa usado na validação do jogo.

\begin{center}
\renewcommand{\arraystretch}{1.1}
\small
\begin{tabular}{@{}>{\raggedright\arraybackslash}p{9.5cm}>{\raggedright\arraybackslash}p{4.5cm}@{}}
\toprule
\textbf{Alteração} & \textbf{Resultado} \\
\midrule
O Cachorro deixa de proteger a Ovelha (R1 sem exceção) & Sem solução \\
Peso máximo do barco igual a 2 & Sem solução \\
Nova regra: Lobo $\times$ Galinha & Sem solução \\
Remoção de R3 (Galinha $\times$ Alface) & Mínimo de 9 travessias \\
Remoção de R4 (Cachorro $\times$ Galinha) & Mínimo de 9 travessias \\
Peso máximo do barco igual a 4 & Mínimo de 5 travessias \\
Lobo com peso 3 e peso máximo do barco igual a 4 & Mínimo de 7 travessias \\
\bottomrule
\end{tabular}
\end{center}

\subsection{Para turmas com formação técnica}

Quem já estudou grafos pode representar cada situação do jogo (quem está em cada
margem e de que lado está o barco) como um círculo, e cada travessia válida como
uma linha entre dois círculos. Das 64 combinações possíveis, 28 são válidas e 26
podem ser alcançadas a partir do início. Uma busca em largura \cite{russell2020}
encontra o caminho mais curto, de 11 travessias. Todo caminho tem número
ímpar de travessias, e por isso 13 é o primeiro limite acima do mínimo.

% ===============================================================
\section{Gabarito (uso do docente)}

Legenda: L = Lobo, O = Ovelha, A = Alface, C = Cachorro, G = Galinha. O
Fazendeiro acompanha todas as travessias.

\textbf{Rodada 1} (sem o Lobo, mínimo de 5 travessias): A e C (E$\rightarrow$D),
viagem vazia (D$\rightarrow$E), O (E$\rightarrow$D), A (D$\rightarrow$E), A e G
(E$\rightarrow$D). A outra solução ótima é a inversa (A e G, A, O, vazia, A e C).

\textbf{Rodada 2} (mínimo de 11 travessias):

\begin{center}
\renewcommand{\arraystretch}{1.0}
\small
\begin{tabular}{@{}rllll@{}}
\toprule
\textbf{N.º} & \textbf{Sentido} & \textbf{Carga} & \textbf{Margem E} & \textbf{Margem D} \\
\midrule
0  & início              & nenhuma & L, O, A, C, G & nenhum        \\
1  & E$\rightarrow$D     & A, G    & L, O, C       & A, G          \\
2  & D$\rightarrow$E     & A       & L, O, C, A    & G             \\
3  & E$\rightarrow$D     & O       & L, C, A       & G, O          \\
4  & D$\rightarrow$E     & nenhuma & L, C, A       & G, O          \\
5  & E$\rightarrow$D     & L       & C, A          & G, O, L       \\
6  & D$\rightarrow$E     & O       & C, A, O       & G, L          \\
7  & E$\rightarrow$D     & A, C    & O             & G, L, A, C    \\
8  & D$\rightarrow$E     & G       & O, G          & L, A, C       \\
9  & E$\rightarrow$D     & O       & G             & L, A, C, O    \\
10 & D$\rightarrow$E     & A       & G, A          & L, C, O       \\
11 & E$\rightarrow$D     & A, G    & nenhum        & L, O, A, C, G \\
\bottomrule
\end{tabular}
\end{center}

\begin{itemize}
  \item Existem exatamente duas soluções com 11 travessias. A segunda coincide
        com a primeira nas travessias 1 a 3 e segue com: G (D$\rightarrow$E),
        A e C (E$\rightarrow$D), O (D$\rightarrow$E), L (E$\rightarrow$D),
        viagem vazia (D$\rightarrow$E), O (E$\rightarrow$D), A
        (D$\rightarrow$E), A e G (E$\rightarrow$D).
  \item O primeiro movimento válido da rodada 2 é único: levar A e G. Qualquer
        outra carga deixa na margem E um par em conflito.
  \item As dicas estão no Apêndice A. Na rodada 2, as dicas 2 e 3 levam ao
        primeiro movimento (A e G) e a dica 5 sugere anotar as margens, o que
        revela ciclos.
\end{itemize}

\textbf{Fluxograma esperado.} A Parte 2 da Ficha, com os quatro quadros
preenchidos, está no Apêndice~C. Os textos dos grupos podem variar, mas a
estrutura deve ser a mesma. O registro das travessias (Parte 1), também preenchido, está no
Apêndice~D.

% ===============================================================
\begingroup\footnotesize\setstretch{1.05}
\begin{thebibliography}{9}

\bibitem{anderson2001}
ANDERSON, L. W.; KRATHWOHL, D. R. (org.). \textbf{A taxonomy for learning,
teaching, and assessing}: a revision of Bloom's taxonomy of educational
objectives. New York: Longman, 2001.

\bibitem{bacich2018}
BACICH, L.; MORAN, J. (org.). \textbf{Metodologias ativas para uma educação
inovadora}: uma abordagem teórico-prática. Porto Alegre: Penso, 2018.

\bibitem{barrows1986}
BARROWS, H. S. A taxonomy of problem-based learning methods.
\textbf{Medical Education}, v. 20, n. 6, p. 481-486, 1986.

\bibitem{bell2009}
BELL, T.; ALEXANDER, J.; FREEMAN, I.; GRIMLEY, M. Computer science unplugged:
school students doing real computing without computers. \textbf{The New Zealand
Journal of Applied Computing and Information Technology}, v. 13, n. 1,
p. 20-29, 2009.

\bibitem{hmelo2004}
HMELO-SILVER, C. E. Problem-based learning: what and how do students learn?
\textbf{Educational Psychology Review}, v. 16, n. 3, p. 235-266, 2004.

\bibitem{russell2020}
RUSSELL, S.; NORVIG, P. \textbf{Artificial intelligence}: a modern approach.
4. ed. Hoboken: Pearson, 2020.

\bibitem{wing2006}
WING, J. M. Computational thinking. \textbf{Communications of the ACM}, v. 49,
n. 3, p. 33-35, 2006.

\end{thebibliography}
\endgroup

% ===============================================================
\appendix
\clearpage
\section{Material desplugado para impressão}

\subsection{Orientações de preparação}

O material foi desenhado para ser impresso e recortado antes da aula. As
ilustrações são originais e foram feitas em TikZ, de modo que este mesmo arquivo
\LaTeX{} gera o conjunto completo.

\begin{center}
\renewcommand{\arraystretch}{1.15}
\small
\begin{tabular}{@{}>{\raggedright\arraybackslash}p{4.0cm}>{\raggedright\arraybackslash}p{3.2cm}>{\raggedright\arraybackslash}p{7.6cm}@{}}
\toprule
\textbf{Item} & \textbf{Quantidade} & \textbf{Observações} \\
\midrule
Tabuleiro (primeira folha de material) & 1 por grupo & Imprimir em A4, orientação paisagem, \textbf{escala 100\%} (sem ``ajustar à página''), pois as peças foram dimensionadas para encaixar nos espaços. \\
Folha de peças e barco (segunda folha de material) & 1 folha para cada 4 grupos & Papel de 180 g/m$^2$ ou colada em papel-cartão. Recortar pelo contorno. \\
Cartas de dica (abaixo) & 1 conjunto por turma & Recortar e entregar somente quando pedidas (regra 8). \\
Regulamento (A.2) & 1 por grupo & Uma página. \\
Ficha de registro (Apêndice B) & 1 por grupo & Três páginas. \\
Fluxograma e ficha de referência (Apêndices C e D) & 1 cópia para o docente & Uso do docente. Podem ser projetados na socialização. \\
Lápis e tesoura & um de cada por grupo & O lápis serve para riscar as gotas de combustível. \\
\bottomrule
\end{tabular}
\end{center}

Exemplo: para 20 estudantes em grupos de 4, imprimem-se 5 tabuleiros, 2 folhas
de peças (a segunda usada pela metade), 1 conjunto de dicas (8 cartas), 5 regulamentos e 5
fichas.

\clearpage
\subsubsection*{Cartas de dica (recortar)}

\noindent\textbf{Rodada 1}\par\smallskip
\noindent
\dica[1]{1}{0,25 ponto}{Há dois primeiros movimentos válidos. Descubram quais são eliminando os outros, um a um, pelas regras.}
\hfill
\dica[1]{2}{0,50 ponto}{Sem o Lobo, os conflitos são Ovelha com Alface, Galinha com Alface e Cachorro com Galinha. Os dois primeiros barcos levam a Alface e mais uma peça.}
\hfill
\dica[1]{3}{0,75 ponto}{A Ovelha atravessa sozinha. Quando ela chega à margem D, a Alface não pode estar lá: a Alface volta no barco. Depois falta levar as duas peças que ficaram.}

\bigskip\noindent\textbf{Rodada 2}\par\smallskip
\noindent
\dica{1}{0,25 ponto}{Há um único primeiro movimento válido. Descubram qual é eliminando os outros, um a um, pelas regras.}
\hfill
\dica{2}{0,25 ponto}{Olhem a margem E depois que o barco sai. Alface com Ovelha ou com Galinha é conflito. Se a Alface ficar, sempre sobra um conflito: ela vai no primeiro barco.}
\hfill
\dica{3}{0,50 ponto}{A Alface vai, mas sozinha não basta: Cachorro e Galinha ficariam juntos. Se o Cachorro for, Lobo e Ovelha ficam sem protetor. Então a Galinha vai: Alface e Galinha.}

\bigskip\noindent
\dica{4}{0,50 ponto}{Nem toda travessia precisa levar carga, e trazer alguém de volta pode ser necessário. Procurem situações em que o Cachorro protege a Ovelha.}
\hfill
\dica{5}{0,75 ponto}{Anotem quem está em cada margem depois de cada travessia. Se uma situação já apareceu antes, vocês estão andando em círculo.}
\hfill
\hspace{5.0cm}


% ---- Tabuleiro: página A4 paisagem em tamanho real ----
\clearpage
\newgeometry{margin=0.9cm,nohead,nofoot}
\begin{landscape}
\thispagestyle{empty}
\noindent\tabuleiro
\end{landscape}

% ---- Peças e barco: quatro kits por folha ----
\clearpage
\newgeometry{margin=1.5cm,nohead,nofoot}
\thispagestyle{empty}
\noindent\paginapecas
\clearpage
\restoregeometry

\subsection{Regulamento para os estudantes}

\begin{tcolorbox}[colback=white,colframe=blue!50!black,coltitle=white,
  colbacktitle=blue!50!black,title={\bfseries Travessia do Rio},
  boxrule=1pt,arc=3mm,fontupper=\small\setstretch{1.12}]
\textbf{Objetivo.} Levar o Lobo, a Ovelha, a Alface, o Cachorro e a Galinha da
margem E para a margem D sem quebrar nenhuma regra e sem esvaziar o tanque de
combustível. O jogo acaba quando as cinco peças estiverem na margem D.

\smallskip
\textbf{Três palavras do jogo.} \textit{Peso} é quanto a peça ocupa no barco.
\textit{Gota} é o combustível gasto, uma por travessia. O jogo não tem pontos: só
a nota do docente usa essa palavra.

\smallskip
\textbf{Papéis.} Cada pessoa assume um papel e todos trocam na rodada 2.
\begin{itemize}[leftmargin=*,itemsep=1pt,topsep=2pt]
  \item \textbf{Operador:} move as peças e o barco.
  \item \textbf{Fiscal:} confere cada travessia contra as regras R1 a R4.
  \item \textbf{Anotador:} preenche a Ficha.
  \item \textbf{Guardião do combustível:} risca uma gota a cada travessia. Em
        grupos de três, o anotador também faz isso.
\end{itemize}

\smallskip
\begin{enumerate}[leftmargin=*,label=\textbf{\arabic*.},itemsep=4pt]
  \item \textbf{Rodadas.} Na rodada 1, o Lobo fica fora do jogo e o tanque tem
        7 gotas. Na rodada 2, todos jogam e o tanque tem 13 gotas.
  \item \textbf{Início.} As peças ficam nos espaços da margem E e o barco, no
        cais E.
  \item \textbf{O barco.} Aguenta peso total 3. O primeiro espaço é do
        Fazendeiro, que pesa 1 e nunca sai do barco. Nos outros dois cabem uma
        peça larga (peso 2), até duas peças quadradas (peso 1 cada) ou nenhuma
        carga.
  \item \textbf{Travessia.} Carreguem o barco e movam-no para o outro cais. A
        carga desembarca na margem de chegada.
  \item \textbf{Conflitos.} Depois de cada travessia, o fiscal examina a margem
        em que o barco \emph{não} está. Se alguma regra R1 a R4 (centro do
        tabuleiro) se aplicar a ela, a travessia é inválida e deve ser
        desfeita. Na margem em que o barco está, e dentro dele, não há
        conflito. As regras valem para qualquer par de peças da margem, mesmo
        que haja outras peças junto: só o Fazendeiro evita o conflito, e a
        única exceção é o Cachorro, que protege a Ovelha do Lobo (R1).
  \item \textbf{Combustível.} Cada travessia gasta uma gota, mesmo se for
        desfeita. Se o tanque esvaziar, recomecem a rodada e anotem a tentativa
        anterior no verso da Ficha.
  \item \textbf{Planejamento.} Simular no papel é permitido e não gasta
        combustível.
  \item \textbf{Dicas.} Há dicas para cada rodada (três na rodada 1 e cinco na
        rodada 2), que podem ser pedidas ao docente. Cada uma reduz a nota do
        grupo: 0,25, 0,50 e 0,75 ponto na rodada 1; 0,25, 0,25, 0,50, 0,50 e
        0,75 ponto na rodada 2, na ordem das cartas.
  \item \textbf{Registro.} Anotem na Ficha as travessias, preencham o
        fluxograma e respondam às questões Q1 a Q3.
\end{enumerate}
\end{tcolorbox}

\clearpage

% ===============================================================
\section{Ficha de registro do grupo}

\noindent\textbf{Grupo:} \rule[-2pt]{3cm}{0.4pt}

\smallskip
\noindent\textbf{Integrantes:} \rule[-2pt]{\dimexpr\linewidth-3.2cm\relax}{0.4pt}

\subsection*{Parte 1: registro das travessias}

Anote quem está em cada margem depois de cada travessia e quantas gotas restam
no tanque.

\begin{tcolorbox}[colback=gray!8,colframe=gray!60,boxrule=0.6pt,arc=1.5mm,
  left=2mm,right=2mm,top=1mm,bottom=1mm,fontupper=\small\setstretch{1.1}]
\textbf{Legenda.} \textbf{L} = Lobo, \textbf{O} = Ovelha, \textbf{A} = Alface,
\textbf{C} = Cachorro, \textbf{G} = Galinha. \textbf{E} = margem esquerda,
\textbf{D} = margem direita. Sentido do barco: E$\rightarrow$D ou
D$\rightarrow$E.
\end{tcolorbox}

\smallskip
\noindent\textbf{Rodada 1} (sem o Lobo; tanque de 7 gotas)

\begin{singlespace}
\begingroup\centering
\renewcommand{\arraystretch}{1.25}
\small
\begin{tabular}{|c|c|p{2.3cm}|p{3.3cm}|p{3.3cm}|c|}
\hline
\textbf{N.º} & \textbf{Sentido} & \textbf{Carga} & \textbf{Margem E} & \textbf{Margem D} & \textbf{Gotas} \\
\hline
0 & & & O, A, C, G & & 7 \\ \hline
1 & & & & & \\ \hline
2 & & & & & \\ \hline
3 & & & & & \\ \hline
4 & & & & & \\ \hline
5 & & & & & \\ \hline
6 & & & & & \\ \hline
7 & & & & & \\ \hline
\end{tabular}
\par\endgroup
\end{singlespace}

\smallskip
\noindent\textbf{Rodada 2} (todos os personagens; tanque de 13 gotas)

\begin{singlespace}
\begingroup\centering
\renewcommand{\arraystretch}{1.25}
\small
\begin{tabular}{|c|c|p{2.3cm}|p{3.3cm}|p{3.3cm}|c|}
\hline
\textbf{N.º} & \textbf{Sentido} & \textbf{Carga} & \textbf{Margem E} & \textbf{Margem D} & \textbf{Gotas} \\
\hline
0 & & & L, O, A, C, G & & 13 \\ \hline
1 & & & & & \\ \hline
2 & & & & & \\ \hline
3 & & & & & \\ \hline
4 & & & & & \\ \hline
5 & & & & & \\ \hline
6 & & & & & \\ \hline
7 & & & & & \\ \hline
8 & & & & & \\ \hline
9 & & & & & \\ \hline
10 & & & & & \\ \hline
11 & & & & & \\ \hline
12 & & & & & \\ \hline
13 & & & & & \\ \hline
\end{tabular}
\par\endgroup
\end{singlespace}

\newpage
\subsection*{Parte 2: a rotina de verificação}

O fluxograma abaixo descreve o que o grupo faz a cada travessia. Complete os
quadros em branco. Cada losango é uma pergunta que se responde com sim ou não;
cada retângulo é uma ação. Escreva como se estivesse explicando para alguém que
nunca jogou.

\begin{center}
\resizebox{!}{0.7\textheight}{\fluxograma
  {\lacunaS}
  {\lacunaS}
  {\lacunaL}
  {\lacunaS}}
\end{center}

\newpage
\subsection*{Parte 3: questões}

\begin{enumerate}[label=\textbf{Q\arabic*.}, leftmargin=1.2cm, itemsep=4.2cm]
  \item Conte uma tentativa que deu errado. O que o grupo percebeu e o que
        mudou na tentativa seguinte?
  \item Compare seu fluxograma com o de outro grupo. O que é igual e o que é
        diferente? Os dois levariam ao mesmo resultado?
  \item Que ideias de Computação apareceram neste jogo sem precisar de um
        computador? Você usaria a atividade com outro público? O que mudaria?
\end{enumerate}

% ===============================================================
\clearpage
\section{Fluxograma de referência (uso do docente)}

A figura reproduz a Parte 2 da Ficha com os quatro quadros preenchidos. Use-a
na socialização para conferir os fluxogramas dos grupos. Ela é uma rotina
válida, não a única possível.

\begin{center}
\resizebox{!}{0.6\textheight}{\fluxograma
  {A carga cabe nos 2 espaços?}
  {Alguma regra (R1 a R4) se aplica à margem sem o barco?}
  {Desfazer a travessia (a gota não volta)}
  {Todas as peças estão na margem D?}}
\end{center}

\begin{tcolorbox}[colback=white,colframe=blue!50!black,boxrule=0.8pt,arc=2mm,
  left=3mm,right=3mm,top=1.5mm,bottom=1.5mm,fontupper=\small\setstretch{1.1}]
\textbf{Ao comparar com o fluxograma de um grupo, confira:}
\begin{itemize}[leftmargin=*,itemsep=1pt,topsep=3pt]
  \item \textbf{a repetição:} todo caminho que não termina volta a ``Escolher a
        carga'', e só a pergunta final, respondida com sim, encerra o jogo;
  \item \textbf{as duas verificações} antes de aceitar a travessia: a da carga
        e a das regras na margem sem o barco;
  \item \textbf{o desfazer:} a travessia inválida é desfeita, mas a gota
        continua gasta;
  \item \textbf{perguntas de sim ou não:} ``Alguém seria devorado ou atacado?''
        serve; ``A jogada foi boa?'' não serve.
\end{itemize}
\end{tcolorbox}

% ===============================================================
\clearpage
\section{Ficha preenchida (uso do docente)}

\newcommand{\pre}[1]{\textcolor{blue!60!black}{#1}}

Registro válido da Parte 1 da Ficha, com o preenchimento em azul. A Parte 2
preenchida está no Apêndice~C. Há outras soluções válidas, de 11 ou 13
travessias, e um grupo pode registrar um caminho diferente e estar correto.
Em cada linha, confira se a margem sem o barco não viola nenhuma regra, se a
carga cabe e se as gotas descem uma por travessia.

\noindent\textbf{Grupo:} \pre{Exemplo} \hfill
\textbf{Integrantes:} \pre{quatro pessoas, um papel cada uma}

\smallskip
\noindent{\small\textbf{Legenda.} \textbf{L} = Lobo, \textbf{O} = Ovelha, \textbf{A} = Alface,
\textbf{C} = Cachorro, \textbf{G} = Galinha. \textbf{E} = margem esquerda,
\textbf{D} = margem direita.}

\smallskip
\noindent\textbf{Rodada 1} (sem o Lobo; tanque de 7 gotas)

\begin{singlespace}
\begingroup\centering
\renewcommand{\arraystretch}{1.3}
\small
\begin{tabular}{|c|c|p{2.3cm}|p{3.3cm}|p{3.3cm}|c|}
\hline
\textbf{N.º} & \textbf{Sentido} & \textbf{Carga} & \textbf{Margem E} & \textbf{Margem D} & \textbf{Gotas} \\
\hline
0 & \pre{início} & \pre{nenhuma} & \pre{O, A, C, G} & \pre{nenhum} & \pre{7} \\ \hline
1 & \pre{E$\rightarrow$D} & \pre{A, C} & \pre{O, G} & \pre{A, C} & \pre{6} \\ \hline
2 & \pre{D$\rightarrow$E} & \pre{nenhuma} & \pre{O, G} & \pre{A, C} & \pre{5} \\ \hline
3 & \pre{E$\rightarrow$D} & \pre{O} & \pre{G} & \pre{O, A, C} & \pre{4} \\ \hline
4 & \pre{D$\rightarrow$E} & \pre{A} & \pre{A, G} & \pre{O, C} & \pre{3} \\ \hline
5 & \pre{E$\rightarrow$D} & \pre{A, G} & \pre{nenhum} & \pre{O, A, C, G} & \pre{2} \\ \hline
6 & & & & & \\ \hline
7 & & & & & \\ \hline
\end{tabular}
\par\endgroup
\end{singlespace}

\smallskip
\noindent\textbf{Rodada 2} (todos os personagens; tanque de 13 gotas)

\begin{singlespace}
\begingroup\centering
\renewcommand{\arraystretch}{1.3}
\small
\begin{tabular}{|c|c|p{2.3cm}|p{3.3cm}|p{3.3cm}|c|}
\hline
\textbf{N.º} & \textbf{Sentido} & \textbf{Carga} & \textbf{Margem E} & \textbf{Margem D} & \textbf{Gotas} \\
\hline
0 & \pre{início} & \pre{nenhuma} & \pre{L, O, A, C, G} & \pre{nenhum} & \pre{13} \\ \hline
1 & \pre{E$\rightarrow$D} & \pre{A, G} & \pre{L, O, C} & \pre{A, G} & \pre{12} \\ \hline
2 & \pre{D$\rightarrow$E} & \pre{A} & \pre{L, O, A, C} & \pre{G} & \pre{11} \\ \hline
3 & \pre{E$\rightarrow$D} & \pre{O} & \pre{L, A, C} & \pre{O, G} & \pre{10} \\ \hline
4 & \pre{D$\rightarrow$E} & \pre{nenhuma} & \pre{L, A, C} & \pre{O, G} & \pre{9} \\ \hline
5 & \pre{E$\rightarrow$D} & \pre{L} & \pre{A, C} & \pre{L, O, G} & \pre{8} \\ \hline
6 & \pre{D$\rightarrow$E} & \pre{O} & \pre{O, A, C} & \pre{L, G} & \pre{7} \\ \hline
7 & \pre{E$\rightarrow$D} & \pre{A, C} & \pre{O} & \pre{L, A, C, G} & \pre{6} \\ \hline
8 & \pre{D$\rightarrow$E} & \pre{G} & \pre{O, G} & \pre{L, A, C} & \pre{5} \\ \hline
9 & \pre{E$\rightarrow$D} & \pre{O} & \pre{G} & \pre{L, O, A, C} & \pre{4} \\ \hline
10 & \pre{D$\rightarrow$E} & \pre{A} & \pre{A, G} & \pre{L, O, C} & \pre{3} \\ \hline
11 & \pre{E$\rightarrow$D} & \pre{A, G} & \pre{nenhum} & \pre{L, O, A, C, G} & \pre{2} \\ \hline
12 & & & & & \\ \hline
13 & & & & & \\ \hline
\end{tabular}
\par\endgroup
\end{singlespace}

\smallskip
\noindent\footnotesize\textbf{Verso da Ficha (exemplo).} Antes de mover as peças,
o grupo simulou no papel levar a Ovelha primeiro. A margem E ficaria com Lobo,
Alface, Cachorro e Galinha, sem o barco: a Galinha bica a Alface (R3) e o
Cachorro ataca a Galinha (R4). Descartaram a ideia sem gastar gota, e o único
primeiro movimento que sobrou foi levar a Alface e a Galinha.
\normalsize

\end{document}